% 问题三方法：风格条件生成、感官画像与双层约束迁移。
\documentclass[tikz,border=0pt]{standalone}
\usepackage[UTF8,fontset=windows]{ctex}
\usepackage{amsmath}
\setmainfont{Times New Roman}
% 方法流程图的共用样式：黑白矩形框、直角箭头、宋体正文、侧边竖排分区标题。
% 坐标单位为毫米，y 轴向下；版心宽 158 mm，侧边标题占 0--6 mm，正文框占 9--158 mm。
\usetikzlibrary{arrows.meta,calc}
% ctex 默认行距系数为 1.3，图内统一取 1。
\newcommand{\fcsize}[2]{\linespread{1}\fontsize{#1}{#2}\selectfont}
\tikzset{
  fc/.style={x=1mm,y=-1mm,>=Latex,line width=0.45pt},
  box/.style={draw,rectangle,anchor=north west,inner sep=1.5mm,outer sep=0pt,align=left,
    font=\songti\fcsize{8.6}{11}},
  % 固定正文高度并顶端对齐，使同一行各框的标题对齐；框外高为 h + 3 mm。
  h/.style={execute at begin node={\begin{minipage}[t][#1][t]{\linewidth}\raggedright},
    execute at end node={\end{minipage}}},
  io/.style={box,dashed},
  key/.style={box,line width=0.9pt},
  edge/.style={->,line width=0.45pt},
  note/.style={font=\songti\fcsize{8}{9.5},fill=white,inner sep=0.6mm},
}
% 框标题
\newcommand{\bt}[1]{{\heiti #1}\\}
% 侧边分区：\sidelabel{上沿 y}{下沿 y}{竖排文字}
\newcommand{\sidelabel}[3]{%
  \draw (0.3,#1) rectangle (6.1,#2);
  \node[align=center,font=\heiti\fcsize{9.5}{11.5},inner sep=0pt]
    at (3.2,{(#1+#2)/2}) {#3};}
% 总线连接：上方节点列表、下方节点列表、总线 y、总线左端 x、右端 x
\newcommand{\bus}[5]{%
  \foreach \u in {#1} {\draw (\u.south) -- (\u.south |- 0,#3);}
  \draw (#4,#3) -- (#5,#3);
  \foreach \l in {#2} {\draw[edge] (\l.north |- 0,#3) -- (\l.north);}}

\begin{document}
\begin{tikzpicture}[fc]
\path[use as bounding box] (0,0) rectangle (158,168);
\sidelabel{0}{28}{输\\入\\与\\表\\示}
\node[io,text width=44mm,h=23mm] (x) at (9,0) {\bt{化学谱与风格}%
250款、22种细分风格；225个单体化合物生成列\\
共同下限视为未检出状态；综合指标另行派生};
\node[box,text width=44mm,h=23mm] (rank) at (60,0) {\bt{逐列秩变换}%
把不同尺度映入潜在高斯空间；并列下限单独处理\\
生成后用经验分位数映回发布尺度};
\node[io,text width=44mm,h=23mm] (up) at (111,0) {\bt{上游模型}%
问题一的纯化学感官预测与等级分布\\
问题二的化合物共表达模块};
\sidelabel{32}{54}{条\\件\\生\\成}
\node[key,text width=95mm,h=18mm] (gen) at (9,32) {\bt{风格条件潜在高斯模型}%
各风格均值向总体均值收缩；风格内共享协方差＝共同因子＋按共表达模块分块的残差\\
在给定风格下抽样，反变换为完整化合物谱};
\node[box,text width=44mm,h=18mm] (tune) at (111,32) {\bt{训练折内选参}%
因子数、均值收缩与协方差结构由留出啤酒的潜空间评分选择};
\bus{x,rank,up}{gen,tune}{29}{26}{140.5}
\sidelabel{59}{91}{生\\成\\核\\验}
\node[box,text width=44mm,h=27mm] (fit) at (9,59) {\bt{保真度}%
分类器区分生成谱与留出真实谱：整体、附风格条件、逐风格\\
比较下限比例及化合物相关结构};
\node[box,text width=44mm,h=27mm] (novel) at (60,59) {\bt{新颖性与可用性}%
生成谱到最近训练啤酒的距离及近复制率\\
用生成谱训练风格分类器，再预测留出真实啤酒};
\node[box,text width=44mm,h=27mm] (ctrl) at (111,59) {\bt{对照生成器}%
无风格条件、逐列独立生成：检验风格和共变作用\\
复制真实啤酒、风格内插值：检验新颖性};
\bus{gen,tune}{fit,novel,ctrl}{55.5}{32.5}{134.5}
\sidelabel{98}{120}{画\\像\\与\\指\\纹}
\node[box,text width=69.5mm,h=18mm] (profile) at (9,98) {\bt{感官画像}%
每种风格生成1000个化合物谱，送入问题一的化学模型，汇总50维预测评分与等级\\
用同一训练折的模型比较留出真实啤酒};
\node[box,text width=69.5mm,h=18mm] (finger) at (85.5,98) {\bt{风味指纹}%
列出相对全体最偏离的化合物和感官维度\\
只描述模型分布，不将关联解释为化合物因果效应};
\bus{fit,novel,ctrl}{profile,finger}{94.5}{45.25}{121.75}
\sidelabel{128}{168}{跨\\风\\格\\迁\\移}
\node[box,text width=32mm,h=24mm] (move) at (9,128) {\bt{1 潜空间平移}%
按源风格到目标风格的均值差移动，保留单款啤酒的风格内偏离};
\node[key,text width=32mm,h=24mm] (project) at (47,128) {\bt{2 化学约束投影}%
9个缺陷化合物：已超阈不得增加，未超阈不得越阈；无阈值不得增加};
\node[key,text width=32mm,h=24mm] (path) at (85,128) {\bt{3 感官约束搜索}%
沿源谱到投影终点的路径，取6个缺陷感官预测等级均不升高的最大可行步长};
\node[box,text width=32mm,h=24mm] (check) at (123,128) {\bt{4 原尺度复查}%
逐列复核两层约束，再评估目标风格程度、改动量与保留的个体特征};
\draw[edge] (move.east) -- (project.west);
\draw[edge] (project.east) -- (path.west);
\draw[edge] (path.east) -- (check.west);
\bus{profile,finger}{move}{124.5}{26}{121.75}
\node[io,text width=146mm,h=9mm] (out) at (9,158) {\bt{交付与对照}%
逐款报告Pils/Lager到Tripel的迁移；汇总全部有序风格对，并与原谱不变、直接截断和检索替代比较。};
\bus{move,project,path,check}{out}{154.5}{26}{140.5}
\end{tikzpicture}
\end{document}
